\section{Primero se elige la representación y después el Transformer}

La sección anterior explicó por qué las arquitecturas tienden a unificarse; esta sección responde la pregunta previa más importante a esa unificación: qué es exactamente lo que debe modelar el Transformer. El video RGB original es continuo y redundante; si se despliega directamente por píxeles, la longitud de la secuencia hace inviable el entrenamiento. Movie Gen usa un Temporal Autoencoder (TAE, autocodificador temporal) para comprimir tanto el espacio como el tiempo en un espacio latent (de variables latentes) y después convierte el latent en patches que se usan como token.

\subsection{Qué variable aleatoria aprende el modelo generativo}

«Aprender la distribución de los videos» exige definir primero la variable aleatoria $x$. Puede tratarse de píxeles, códigos discretos, un latent continuo o una representación multiescala. Las dos diapositivas siguientes comparan el texto con los medios audiovisuales: los token del lenguaje se prestan casi de forma natural al modelado de secuencias, mientras que el video necesita eliminar antes una gran cantidad de redundancia predecible.

\lecturefigure{slide-177-architecture-questions.jpg}{Las tres preguntas de la arquitectura de Movie Gen: representación, objetivo de aprendizaje y backbone}{00:11:19}

Estas tres preguntas organizan toda la clase: la representation determina la longitud en token y la pérdida de información; el learning objective determina qué predice la red; la architecture determina cómo entran las condiciones y cómo escala el cómputo. Las tres se restringen entre sí, por lo que no tiene sentido comparar de forma aislada «qué bloque de Transformer es más potente».

\lecturefigure{slide-178-representation-definition.jpg}{El punto de partida del aprendizaje de representaciones: para aprender $P(x)$ hay que definir primero $x$}{00:11:43}

El texto puede convertirse en token discretos mediante un tokenizer; el video, en cambio, es un tensor continuo de $T_0\times3\times H_0\times W_0$. Si cada valor RGB se tratara como un token, la dimensión de los datos y la longitud de la secuencia superarían con mucho las del texto, y además los píxeles y los cuadros contiguos son altamente redundantes.

\lecturefigure{slide-181-text-vs-media-representation.jpg}{La representación del texto es compacta y discreta; la de los medios audiovisuales es continua y redundante}{00:13:36}

Al leer la figura, conviene comparar primero la «densidad semántica»: un token de palabra puede resumir un concepto complejo, mientras que un píxel solo expresa un color local. El objetivo de la representación de medios no es comprimir sin pérdida todos los detalles, sino conservar, con un error de reconstrucción aceptable, los objetos, las texturas, el movimiento y la estructura temporal que necesita la tarea de generación.

\begin{importantbox}{Elegir la representación es repartir el presupuesto de cómputo}
Una compresión más fuerte acorta la secuencia y reduce el costo del Transformer, pero puede perder objetos pequeños, movimientos rápidos y texturas finas; una compresión más débil conserva información, pero traslada el costo a la atención, las activaciones y la comunicación. La representación no es un preprocesamiento de datos, sino la primera capa de la arquitectura del sistema generativo.
\end{importantbox}

\subsection{Por qué no generar píxeles directamente}

La sección anterior mostró la redundancia de los medios; esta convierte esa redundancia en número de token. Sea el video de entrada
\[
V\in\mathbb R^{T_0\times 3\times H_0\times W_0},
\]
donde $T_0$ es el número de cuadros, $H_0,W_0$ son el alto y el ancho en píxeles, y el número de canales es 3. Si se desplegara directamente cada posición, el tamaño de la secuencia sería aproximadamente $T_0H_0W_0$, lo cual es inviable incluso ignorando los canales de color.

\lecturefigure{slide-182-latent-representation.jpg}{Modelar directamente los píxeles es conceptualmente sencillo, pero el video de alta resolución requiere un sobremuestreo complejo o una cantidad enorme de cómputo}{00:15:52}

Movie Gen elige un latent continuo en lugar de los píxeles directos. El encoder del TAE transforma la entrada en
\[
X=E(V)\in\mathbb R^{T\times C\times H\times W},
\qquad
\frac{T_0}{T}=\frac{H_0}{H}=\frac{W_0}{W}=8,
\]
donde $C=16$ es el número de latent channels. Cada uno de los tres ejes se comprime 8 veces, de modo que el número de posiciones espaciotemporales se reduce aproximadamente $8^3=512$ veces.

\lecturefigure{slide-185-temporal-autoencoder.jpg}{El TAE comprime el video RGB en un latent espaciotemporal continuo y el decoder recupera los píxeles}{00:19:17}

Para un video de 256 cuadros y $768\times768$, tras la compresión quedan unas $32\times96\times96$ posiciones espaciotemporales; si cada posición correspondiera a un token, entonces
\[
N=\frac{256}{8}\cdot\frac{768}{8}\cdot\frac{768}{8}
=294{,}912.
\]
El context de 73K del artículo incluye además un patchify adicional, por lo que el número de token disminuye todavía más. Este ejemplo muestra que una «compresión de 8 veces» debe precisar sobre qué ejes actúa; no debe interpretarse como si la compresión total fuera de solo 8 veces.

\begin{knowledgebox}{Cómo el TAE admite a la vez imágenes y videos}
El TAE añade una convolución temporal unidimensional después de la convolución espacial bidimensional, y una atención temporal después de la atención espacial. Una imagen se trata como un video de un solo cuadro; mediante una codificación de longitud variable y el descarte de los cuadros decodificados sobrantes, el mismo encoder/decoder puede procesar imágenes y videos de cualquier longitud.
\end{knowledgebox}

\begin{warningbox}{Una tasa de compresión alta no conserva la información sin costo}
Una pérdida de reconstrucción baja no significa que se conserven todas las propiedades relevantes para la generación. Los movimientos rápidos, las líneas finas, el texto, los rostros pequeños y las texturas parpadeantes pueden promediarse en el latent. El TAE debe evaluarse a la vez en reconstrucción espacial, consistencia temporal y calidad de la generación posterior; no basta con reportar el PSNR de cuadros individuales.
\end{warningbox}

\subsection{La pérdida del TAE y los límites de la adversarial loss}

El TAE pertenece a la línea de los VAE, pero el artículo añade además objetivos como la pérdida perceptual, la del discriminator y la outlier penalty. Para la reconstrucción $\hat V=D(E(V))$, el objetivo global puede escribirse como
\[
\mathcal L_{\mathrm{TAE}}
=\lambda_{\mathrm{rec}}\mathcal L_{\mathrm{rec}}
+\lambda_{\mathrm{perc}}\mathcal L_{\mathrm{perc}}
+\lambda_{\mathrm{adv}}\mathcal L_{\mathrm{adv}}
+\lambda_{\mathrm{KL}}\mathcal L_{\mathrm{KL}}
+\lambda_{\mathrm{OPL}}\mathcal L_{\mathrm{OPL}}.
\]
Cada $\lambda$ es un peso de la pérdida; el término de reconstrucción restringe el contenido, el término perceptual restringe la similitud visual semántica, el término adversarial favorece el realismo, el término KL regulariza el latent y el OPL penaliza los latent dots con norma anormalmente alta.

\begin{knowledgebox}{Primera aparición: adversarial loss}
Aquí la adversarial loss proviene de un discriminador que se usa para mejorar el realismo de la decodificación y la capacidad de compresión del autoencoder; no significa que el generador principal de Movie Gen sea una GAN. El generador principal se sigue entrenando con Flow Matching. En la sesión de Q\&A, el docente enfatiza que añadir este tipo de pérdida permite que el decoder no tenga que copiar la entrada píxel por píxel, lo que hace posible una compresión mayor.
\end{knowledgebox}

\teachervoice{01:07:20--01:09:23, el docente explica que la pérdida del discriminador del TAE proviene de la línea de VQGAN: la restricción L1 a nivel de píxel es demasiado rígida, mientras que la adversarial loss da a la reconstrucción más libertad perceptual y aporta una ganancia adicional de compresión; esto es algo distinto de «generar video con una GAN».}

\begin{lstlisting}[language=Python,caption={Cálculo mínimo de la compresión de video y del token accounting}]
latent = temporal_autoencoder.encode(video)  # [T/8, C=16, H/8, W/8]
tokens = patchify(latent, patch_t, patch_h, patch_w)
sequence_length = tokens.shape[0]
attention_pairs = sequence_length ** 2
report(sequence_length, attention_pairs, reconstruction_error)
\end{lstlisting}

\subsection{Resumen de la sección}

Movie Gen no hace que el Transformer modele directamente los píxeles: primero usa el TAE para comprimir 8 veces cada uno de los ejes de tiempo, alto y ancho, y después aplica patchify para obtener token. La compresión vuelve factible el cómputo, pero también introduce la reconstrucción y un cuello de botella de información. Que el TAE use adversarial loss no significa que el generador principal sea una GAN; el objetivo de aprendizaje del backbone generativo es el Flow Matching de la sección siguiente.
